% Références : celles de l'application (code_source/cora/app/modules/10_base.js, vérifiées sur PubMed le 13 et le
% 16/10/2026) et des ouvrages ou articles de méthode utilisés par le code.
\newcommand{\doi}[1]{\href{https://doi.org/#1}{doi:\nolinkurl{#1}}}
\cleardoublepage
\phantomsection
\addcontentsline{toc}{chapter}{Bibliographie}
\markboth{Bibliographie}{}
\begin{thebibliography}{99}
\small\setlength{\parskip}{0.15em}\setlength{\itemsep}{0.2em}

\bibitem{paley2002} Paley D. \emph{Principles of Deformity Correction}. Springer, 2002. \doi{10.1007/978-3-642-59373-4}

\bibitem{tomlinson2007} Tomlinson J, Fox D, Cook JL, Keller GG. Measurement of femoral angles in four dog breeds.
\emph{Vet Surg} 2007;36:593--598. \doi{10.1111/j.1532-950X.2007.00309.x}

\bibitem{dismukes2007} Dismukes DI, Tomlinson JL, Fox DB, Cook JL, Song KJ. Radiographic measurement of the proximal
and distal mechanical joint angles in the canine tibia. \emph{Vet Surg} 2007;36:699--704.
\doi{10.1111/j.1532-950X.2007.00323.x}

\bibitem{dismukes2008s} Dismukes DI, Fox DB, Tomlinson JL, Essman SC. Radiographic measurement of canine tibial
angles in the sagittal plane. \emph{Vet Surg} 2008;37:300--305. \doi{10.1111/j.1532-950X.2008.00381.x}

\bibitem{kara2018} Kara ME, Sevil-Kilimci F, Dilek ÖG, Onar V. Proximal and distal alignment of normal canine femurs:
a morphometric analysis. \emph{Ann Anat} 2018;217:125--128. \doi{10.1016/j.aanat.2018.02.006}

\bibitem{olsen2016} Olsen AM, Vezzoni L, Ferretti A, Palmer RH, Vezzoni A, Duerr F. Hemiepiphysiodesis for the
correction of proximal tibial valgus in growing dogs. \emph{Vet Comp Orthop Traumatol} 2016;29(4):330--337.
\doi{10.3415/VCOT-15-12-0204}

\bibitem{vezzoni2022} Vezzoni L, Forzisi I, Ferretti A, Vezzoni A. Hemiepiphysiodesis for the correction of distal
femoral valgus in growing dogs. \emph{Vet Comp Orthop Traumatol} 2022;35(2):112--118. \doi{10.1055/s-0041-1736381}

\bibitem{vezzoni2020} Vezzoni L, Bazzo S, Boiocchi S, Vezzoni A. Modified TPLO with double cut and medial crescentic
closing wedge osteotomy to treat dogs with CCL rupture and tibial valgus deformity. \emph{Vet Comp Orthop Traumatol}
2020;33(1):59--65. \doi{10.1055/s-0039-1700565}

\bibitem{guenego2021} Guénégo L, Vezzoni A, Vezzoni L. Comparison of tibial anatomical-mechanical axis angles and
patellar positions between TPLO and AMA-based CCWO and clinically normal Labrador retrievers. \emph{BMC Vet Res}
2021;17:368. \doi{10.1186/s12917-021-03094-3}

\bibitem{guenego2025} Guénégo L, Vezzoni A, Vezzoni L. Modified cranial closing wedge osteotomy (AMA-based CCWO) in
204 dogs over 50 kg. \emph{Vet Comp Orthop Traumatol} 2025;38(3):111--118. \doi{10.1055/s-0044-1790511}

\bibitem{wilke2002} Wilke VL, Conzemius MG, Besancon MF, Evans RB, Ritter M. Comparison of tibial plateau angle
between clinically normal Greyhounds and Labrador Retrievers with and without rupture of the cranial cruciate ligament.
\emph{J Am Vet Med Assoc} 2002;221(10):1426--1429. \doi{10.2460/javma.2002.221.1426}

\bibitem{imaios} IMAIOS vet-Anatomy. \emph{Valgus tibial} ; \emph{Angle mécanique médial proximal du tibia} ;
\emph{Angle de varus fémoral}. \url{https://www.imaios.com/fr/vet-anatomy}

\bibitem{bruhschwein2023r} Brühschwein A, Burg-Personnaz J, Zöllner M, Reese S, Meyer-Lindenberg A. Comparison of
CT-measured angles of pelvic limbs without patellar luxation of six canine breeds. \emph{Front Vet Sci}
2023;10:1194167. \doi{10.3389/fvets.2023.1194167}

\bibitem{bruhschwein2023f} Brühschwein A, Schmitz B, Zöllner M, Reese S, Meyer-Lindenberg A. CT angular measurements
using a bone-centered 3D coordinate system are accurate in a femoral torsional deformity model. \emph{Front Vet Sci}
2023;10:1019216. \doi{10.3389/fvets.2023.1019216}

\bibitem{bruhschwein2023t} Brühschwein A, Schmitz B, Zöllner M, Reese S, Meyer-Lindenberg A. Three-dimensional
computed tomographic angular measurements of the canine tibia using a bone-centered coordinate system. \emph{Front Vet
Sci} 2023;10:1154144. \doi{10.3389/fvets.2023.1154144}

\bibitem{alaiyan2021} Al Aiyan A, Richardson K, Manchi G, Ginja M, Brunnberg L. Measurement of the femoral anteversion
angle in medium and large dog breeds using computed tomography. \emph{Front Vet Sci} 2021;8:540406.
\doi{10.3389/fvets.2021.540406}

\bibitem{longo2022} Longo F, Nicetto T, Knell SC, Evans RB, Isola M, Pozzi A. Three-dimensional volume rendering
planning, surgical treatment, and clinical outcomes for femoral and tibial detorsional osteotomies in dogs. \emph{Vet
Surg} 2022;51(7):1126--1141. \doi{10.1111/vsu.13882}

\bibitem{dudley2006} Dudley RM, Kowaleski MP, Drost WT, Dyce J. Radiographic and computed tomographic determination of
femoral varus and torsion in the dog. \emph{Vet Radiol Ultrasound} 2006;47(6):546--552.
\doi{10.1111/j.1740-8261.2006.00184.x}

\bibitem{apelt2005} Apelt D, Kowaleski MP, Dyce J. Comparison of computed tomographic and standard radiographic
determination of tibial torsion in the dog. \emph{Vet Surg} 2005;34(5):457--462. \doi{10.1111/j.1551-2916.2005.00069.x}

\bibitem{aper2005} Aper R, Kowaleski MP, Apelt D, Drost WT, Dyce J. Computed tomographic determination of tibial
torsion in the dog. \emph{Vet Radiol Ultrasound} 2005;46(3):187--191. \doi{10.1111/j.1740-8261.2005.00048.x}

\bibitem{mostafa2014} Mostafa AA, Griffon DJ, Thomas MW, Constable PD. Radiographic evaluation of femoral torsion and
correlation with computed tomographic techniques in Labrador Retrievers. \emph{Vet Surg} 2014;43(5):534--541.
\doi{10.1111/j.1532-950X.2014.12096.x}

\bibitem{miles2016} Miles JE. Femoral rotation unpredictably affects radiographic anatomical lateral distal femoral
angle measurements. \emph{Vet Comp Orthop Traumatol} 2016;29(2):156--159. \doi{10.3415/VCOT-15-06-0107}

\bibitem{ivanov2024} Ivanov DV, Welby JP, Khanna A, et al. Evaluation of intraoperative fluoroscopic techniques to
estimate femoral rotation: a cadaveric study. \emph{J Orthop Trauma} 2024;38(5):279--284.
\doi{10.1097/BOT.0000000000002790}

\bibitem{krettek1998} Krettek C, Miclau T, Grün O, Schandelmaier P, Tscherne H. Intraoperative control of axes,
rotation and length in femoral and tibial fractures. \emph{Injury} 1998;29 Suppl 3:C29--39.
\doi{10.1016/s0020-1383(98)95006-9}

\bibitem{fang2015} Fang C, Gibson W, Lau TW, et al. Important tips and numbers on using the cortical step and diameter
difference sign in assessing femoral rotation: should we abandon the technique? \emph{Injury} 2015;46(7):1393--1399.
\doi{10.1016/j.injury.2015.04.009}

\bibitem{kim2017} Kim SY, Snowdon KA, DeCamp CE. Single oblique osteotomy for correction of antebrachial angular and
torsional deformities in a dog. \emph{J Am Vet Med Assoc} 2017;251(3):333--339. \doi{10.2460/javma.251.3.333}

\bibitem{qatu2024} Qatu M, Al-Jamal M, Mouzaihem H. Rigorous mathematical derivation of inclination angle for single
cut rotational osteotomy for tibia deformity correction. \emph{Orthop Rev} 2024;16:120302. \doi{10.52965/001c.120302}

\bibitem{petazzoni2012} Petazzoni M, Nicetto T, Vezzoni A, Piras A, Palmer R. Treatment of pes varus using locking
plate fixation in seven Dachshund dogs. \emph{Vet Comp Orthop Traumatol} 2012;25(3):231--238.
\doi{10.3415/VCOT-11-03-0035}

\bibitem{brower2017} Brower BE, Kowaleski MP, Peruski AM, et al. Distal femoral lateral closing wedge osteotomy as a
component of comprehensive treatment of medial patellar luxation and distal femoral varus in dogs. \emph{Vet Comp
Orthop Traumatol} 2017;30(1):20--27. \doi{10.3415/VCOT-16-07-0103}

\bibitem{noyes2005} Noyes FR, Goebel SX, West J. Opening wedge tibial osteotomy: the 3-triangle method to correct axial
alignment and tibial slope. \emph{Am J Sports Med} 2005;33(3):378--387. \doi{10.1177/0363546504269034}

\bibitem{ducat2012} Ducat A, Sariali E, Lebel B, et al. Posterior tibial slope changes after opening- and
closing-wedge high tibial osteotomy. \emph{Orthop Traumatol Surg Res} 2012;98(1):68--74. \doi{10.1016/j.otsr.2011.08.013}

\bibitem{fox2006} Fox DB, Tomlinson JL, Cook JL, Breshears LM. Principles of uniapical and biapical radial deformity
correction using dome osteotomies and the center of rotation of angulation methodology in dogs. \emph{Vet Surg}
2006;35(1):67--77. \doi{10.1111/j.1532-950X.2005.00114.x}

\bibitem{worth2019} Worth AJ, Crosse KR, Kersley A. Computer-assisted surgery using 3D printed saw guides for acute
correction of antebrachial angular limb deformities in dogs. \emph{Vet Comp Orthop Traumatol} 2019;32(3):241--249.
\doi{10.1055/s-0039-1678701}

\bibitem{hall2019} Hall EL, Baines S, Bilmont A, Oxley B. Accuracy of patient-specific three-dimensional-printed
osteotomy and reduction guides for distal femoral osteotomy in dogs with medial patella luxation. \emph{Vet Surg}
2019;48(4):584--591. \doi{10.1111/vsu.13126}

\bibitem{carwardine2021} Carwardine DR, Gosling MJ, Burton NJ, O'Malley FL, Parsons KJ. Three-dimensional-printed
patient-specific osteotomy guides, repositioning guides and titanium plates for acute correction of antebrachial limb
deformities in dogs. \emph{Vet Comp Orthop Traumatol} 2021;34(1):43--52. \doi{10.1055/s-0040-1709702}

\bibitem{dearmond2022} De Armond CC, Lewis DD, Kim SE, Biedrzycki AH. Accuracy of virtual surgical planning and custom
three-dimensionally printed osteotomy and reduction guides for acute uni- and biapical correction of antebrachial
deformities in dogs. \emph{J Am Vet Med Assoc} 2022;260(13):1--9. \doi{10.2460/javma.21.09.0419}

\bibitem{jeon2023} Jeon JW, Kang KW, Kim WK, Jung C, Kang BJ. Three-dimensional-printed patient-specific guides for
tibial deformity correction in small-breed dogs. \emph{Am J Vet Res} 2023;84(12). \doi{10.2460/ajvr.23.06.0146}

\bibitem{bright2024} Bright S, Schofield I, Oxley B. Surgical correction of pes varus deformity in dachshunds using
three-dimensional-printed patient-specific guide system: nine tibiae in seven cases (2018--2022). \emph{J Small Anim
Pract} 2024;65(12):894--905. \doi{10.1111/jsap.13789}

\bibitem{hounsfield1973} Hounsfield GN. Computerized transverse axial scanning (tomography). Part 1: Description of
system. \emph{Br J Radiol} 1973;46(552):1016--1022.

\bibitem{dicom} NEMA. \emph{Digital Imaging and Communications in Medicine (DICOM) Standard}, partie PS3.3 :
définitions des objets. \url{https://www.dicomstandard.org}

\bibitem{fedorov2012} Fedorov A, Beichel R, Kalpathy-Cramer J, et al. 3D Slicer as an image computing platform for
the Quantitative Imaging Network. \emph{Magn Reson Imaging} 2012;30(9):1323--1341.

\bibitem{meyer1994} Meyer F. Topographic distance and watershed lines. \emph{Signal Processing} 1994;38(1):113--125.

\bibitem{lorensen1987} Lorensen WE, Cline HE. Marching cubes: a high resolution 3D surface construction algorithm.
\emph{Computer Graphics (SIGGRAPH '87)} 1987;21(4):163--169.

\bibitem{schroeder2015} Schroeder W, Maynard R, Geveci B. Flying edges: a high-performance scalable isocontouring
algorithm. \emph{IEEE Symposium on Large Data Analysis and Visualization (LDAV)}, 2015.

\bibitem{taubin1995} Taubin G. A signal processing approach to fair surface design. \emph{Proceedings of SIGGRAPH
'95}, 1995:351--358.

\bibitem{garland1997} Garland M, Heckbert PS. Surface simplification using quadric error metrics. \emph{Proceedings of
SIGGRAPH '97}, 1997:209--216.

\bibitem{fischler1981} Fischler MA, Bolles RC. Random sample consensus: a paradigm for model fitting with applications
to image analysis and automated cartography. \emph{Commun ACM} 1981;24(6):381--395.

\bibitem{maurer2003} Maurer CR, Qi R, Raghavan V. A linear time algorithm for computing exact Euclidean distance
transforms of binary images in arbitrary dimensions. \emph{IEEE Trans Pattern Anal Mach Intell} 2003;25(2):265--270.

\bibitem{ericson2005} Ericson C. \emph{Real-Time Collision Detection}. Morgan Kaufmann, 2005 (§ 5.1.5 : point le plus
proche sur un triangle).

\bibitem{quilez2013} Quílez I. \emph{Smooth minimum} (article en ligne). \url{https://iquilezles.org/articles/smin/}

\bibitem{rodrigues1840} Rodrigues O. Des lois géométriques qui régissent les déplacements d'un système solide dans
l'espace. \emph{Journal de mathématiques pures et appliquées} 1840;5:380--440.

\bibitem{press2007} Press WH, Teukolsky SA, Vetterling WT, Flannery BP. \emph{Numerical Recipes: The Art of
Scientific Computing}, 3\textsuperscript{e} éd. Cambridge University Press, 2007.

\bibitem{harris2020} Harris CR, Millman KJ, van der Walt SJ, et al. Array programming with NumPy. \emph{Nature}
2020;585:357--362.

\bibitem{virtanen2020} Virtanen P, Gommers R, Oliphant TE, et al. SciPy 1.0: fundamental algorithms for scientific
computing in Python. \emph{Nat Methods} 2020;17:261--272.

\bibitem{lowekamp2013} Lowekamp BC, Chen DT, Ibáñez L, Blezek D. The design of SimpleITK. \emph{Front Neuroinform}
2013;7:45.

\bibitem{schroeder2006} Schroeder W, Martin K, Lorensen B. \emph{The Visualization Toolkit: An Object-Oriented
Approach to 3D Graphics}, 4\textsuperscript{e} éd. Kitware, 2006.

\bibitem{manifold} Lalish E, et al. \emph{Manifold}, bibliothèque de géométrie aux opérations booléennes robustes
(logiciel libre). \url{https://github.com/elalish/manifold}

\end{thebibliography}
